\subsection{导子的例子}

例 1. *设 A 是 R 到 R 的可微映射构成的 R-代数，并设 \(x_0\) 是 R 的一个点；R 可视为一个 A-模，其外运算为 \((f, a) \mapsto f(x_0)a\) 。则映射 \(f \mapsto \mathrm{D}f(x_0)\) 是一个导子，因为（实变函数，I，§ 1，第 3 号）

\[
(\mathbf {D} (f g)) (x _ {0}) = (\mathbf {D} f (x _ {0})) g (x _ {0}) + f (x _ {0}) (\mathbf {D} g (x _ {0})). *
\]

例 2. *设 X 是一个 C∞ 类可微流形，并设 A 为 X 上微分形式的分次 R-代数。将每个微分形式 ω 映到其外微分 dω 的映射是次数为 +1 的反导子（可微与解析流形，R，§ 8）。*

例 3. 设 A 是一个结合 K-代数。对所有 \(a \in \mathbf{A}\) ，映射 \(x \mapsto ax - xa\) 是代数 A 的一个导子（参见第 6 号）。

例 4. 设 M 是一个 K-模，A 是带通常分次的外代数 \(\bigwedge (\mathbf{M}^{*})\) （§ 7，第 1 号）。将在 § 11 第 9 号中看到，对所有 \(x\in \mathbf{M}\) ，右内积 \(i(x)\) 是 A 的一个次数为 \(-1\) 的反导子。*

例 5. 回到第 2 号定义 1 的一般情形，设 \(\overline{\mathbf{K}}\) 是另一个交换环，且 \(\rho: \mathbf{K} \to \overline{\mathbf{K}}\) 是一个环同态；设 \(\overline{\mathbf{A}}, \overline{\mathbf{A}'}\) ， \(\overline{\mathbf{A}''}\) ， \(\overline{\mathbf{B}}, \overline{\mathbf{B}'}\) ， \(\overline{\mathbf{B}''}\) 表示分次 \(\overline{\mathbf{K}}\) -模，它们分别由 \(\mathbf{A}, \mathbf{A}', \mathbf{A}''\) ， \(\mathbf{B}, \mathbf{B}'\) ， \(\mathbf{B}''\) 通过将标量环扩张到 \(\overline{\mathbf{K}}\) 而得到（II，§ 11，第 5 号）；我们由 \(\mu\) ， \(\lambda_1\) 和 \(\lambda_2\) 导出 \(\overline{\mathbf{K}}\) -双线性映射

\[
\bar {\mu}: \overline {{{A}}} \times \overline {{{A}}} ^ {\prime} \rightarrow \overline {{{A}}} ^ {\prime \prime}, \quad \bar {\lambda} _ {1}: \bar {B} \times \overline {{{A}}} ^ {\prime} \rightarrow \bar {B} ^ {\prime \prime}, \quad \bar {\lambda} _ {2}: \overline {{{A}}} \times \bar {B} ^ {\prime} \rightarrow \bar {B} ^ {\prime \prime}
\]

作 \(l_{\bar{K}}\) 与 \(\mu\) ， \(\lambda_{1}\) 和 \(\lambda_{2}\) 相应的 K-线性映射的张量积（II, § 5, 第 1 号）。那么，若 d 是一个 \(\varepsilon\) -导子（次数为 \(\delta\) ，从 \((\mathbf{A}, \mathbf{A}', \mathbf{A}^{\prime\prime})\) 到 \((\mathbf{B}, \mathbf{B}', \mathbf{B}^{\prime\prime})\) ），则映射 \(\bar{d} = d \otimes l_{\bar{K}}\) （从 \(\overline{A} \oplus \overline{A}' \oplus \overline{A}''\) 到 \(\overline{B} \oplus \overline{B}' \oplus \overline{B}''\) ）是一个 \(\varepsilon\) -导子（次数为 \(\delta\) ，从 \((\overline{A}, \overline{A}', \overline{A}^{\prime\prime})\) 到 \((\overline{B}, \overline{B}', \overline{B}'' )\) ）。

例 6. 设 A 为类型 Z 的分次 K-代数；通过如下方式定义一个次数为 0 的分次 K-线性映射 \(d: A \to A\) ：对 \(x_{n} \in \mathbf{A}_{n}(n \in \mathbf{Z})\) ，取 \(d(x_{n}) = nx_{n}\) 。该映射是 A 的一个导子，因为对 \(x_{p} \in A_{p}, x_{q} \in A_{q}\) ，

\[
d (x _ {p} x _ {q}) = (p + q) x _ {p} x _ {q} = d (x _ {p}) x _ {q} + x _ {p} d (x _ {q}).
\]

